\chapter*{Introducción}\label{chap:intro}
\addcontentsline{toc}{chapter}{Introducción}

La preservación de documentos físicos constituye una necesidad de las instituciones culturales y científicas. Archivos institucionales, fondos bibliográficos, expedientes judiciales y acervos periodísticos acumulan millones de páginas cuyo contenido no está indexado ni procesado por ningún sistema informático~\cite{muehlberger2019transkribus}. La conversión de esos documentos en texto procesable es el punto de partida para su indexación, búsqueda y análisis a escala. Sin esa conversión, el contenido queda fuera del alcance de las herramientas de procesamiento de lenguaje natural, la minería de datos y los sistemas de recuperación de información. La democratización del patrimonio documental de la humanidad depende del avance de las tecnologías de digitalización~\cite{unesco2003charter}.

El reconocimiento óptico de caracteres, denominado \textit{Optical Character Recognition} en la literatura anglosajona, OCR por sus siglas en inglés, es la tecnología que convierte imágenes de documentos en texto digital procesable, y constituye el componente central de las iniciativas de digitalización documental~\cite{subramani2021survey}. El Instituto Cubano de Investigación Cultural Juan Marinello constituye un ejemplo de dicha situación: su acervo documental, en gran parte mecanografiado e impreso en español, no cuenta con transcripciones digitales que permitan su consulta, indexación o análisis computacional.

El OCR para documentos impresos en ese contexto enfrenta un problema de datos. A diferencia del inglés, que cuenta con conjuntos de datos sintéticos con varios millones de muestras~\cite{jaderberg2014synthetic}, el español carece de recursos equivalentes que cubran sus signos tipográficos, como lo son la raya de diálogo, las comillas angulares, los signos de apertura de interrogación y exclamación, la eñe y las vocales con tilde. La distribución de caracteres en español sigue la ley de Zipf~\cite{zipf1949human}, de modo que algunas letras aparecen con frecuencia muy superior a otras, lo cual introduce un desequilibrio de clases que los sistemas de reconocimiento deben compensar.

Varios trabajos abordan problemas relacionados con la digitalización de texto impreso. Shi et al.~\cite{shi2016crnn} propone la arquitectura de red neuronal convolucional recurrente, denominada \textit{Convolutional Recurrent Neural Network} en inglés y conocida por sus siglas CRNN, que combina una red neuronal convolucional con una red neuronal recurrente bidireccional, se entrena con la función de pérdida \textit{Connectionist Temporal Classification}, conocida por sus siglas CTC, y transcribe líneas de texto completas sin segmentación previa de caracteres. Breuel et al.~\cite{breuel2013ocropus} demostró que una red bidireccional LSTM, sin modelo de lenguaje ni postprocesado, alcanza tasas de error de carácter inferiores al~2\,\% sobre texto impreso en inglés y sobre la degradación artificial de la tipografía Fraktur. El resultado estableció la viabilidad de las arquitecturas LSTM para el reconocimiento de líneas impresas, y abrió la puerta al entrenamiento sobre \textit{corpus} con artefactos sintéticos que reproducen las condiciones visuales de los documentos objetivo. Buda et al.~\cite{buda2018imbalance} caracterizó el sesgo que introduce el desequilibrio de frecuencias en clasificadores basados en redes convolucionales, y Chawla et al.~\cite{chawla2002smote} propuso técnicas generales de sobremuestreo para clases minoritarias, aplicables al entrenamiento de sistemas OCR cuando el \textit{corpus} presenta caracteres infrecuentes. En el plano de los sistemas de acceso a documentos digitalizados, Muehlberger et al.~\cite{muehlberger2019transkribus} describe la plataforma Transkribus, que integra modelos de reconocimiento de texto con herramientas de gestión, búsqueda y revisión de transcripciones.

El problema que aborda el presente trabajo es el desarrollo de un sistema OCR para texto impreso en español que pueda integrarse en una aplicación web de digitalización documental. Lo que distingue este trabajo de los anteriores es la combinación de tres elementos: un \textit{corpus} sintético anclado en las frecuencias estadísticas del español para contrarrestar el desequilibrio de clases, un modelo de reconocimiento de secuencias que aprende sobre ese \textit{corpus}, y una aplicación web que integra el sistema con funcionalidades de gestión documental y revisión colaborativa con roles diferenciados.

El objetivo general de este trabajo es diseñar e implementar un sistema de reconocimiento óptico de caracteres para documentos impresos en lengua española e integrarlo en una aplicación web para la digitalización, visualización y gestión de dichos documentos.

Para alcanzar este objetivo, se plantean los siguientes objetivos específicos:

\begin{enumerate}
	\item Realizar una revisión del estado del arte en sistemas OCR, con atención a las arquitecturas de aprendizaje profundo, a los conjuntos de datos disponibles para texto impreso en español y a las plataformas de visualización y gestión de documentos digitalizados.
	
	\item Diseñar y generar un \textit{corpus} sintético de pares imagen--transcripción que cubra los símbolos propios del español impreso, corrija el desequilibrio de frecuencias de caracteres y sea representativo de la diversidad tipográfica de los documentos objetivo, e incluir la validación estadística del \textit{corpus} resultante.
	
	\item Implementar el \textit{pipeline} de preprocesamiento de documentos, que comprende la binarización adaptativa, la corrección de inclinación y la segmentación de líneas de texto, y que produce las muestras que el modelo recibe para el entrenamiento y la inferencia.
	
	\item Diseñar, entrenar y evaluar los resultados de un modelo de reconocimiento de secuencias con arquitectura CRNN, red recurrente bidireccional y función de pérdida CTC. Además, se incluye implementar la decodificación por búsqueda en haz con normalización de longitud y modelo de lenguaje en español para la inferencia.
	
	\item Evaluar el rendimiento del sistema mediante las métricas estándar del campo. El protocolo incluye: tasa de error de carácter, tasa de error de palabra y exactitud por línea, junto con intervalos de confianza al~95\,\% que el procedimiento \textit{bootstrap} produce. El protocolo realiza también pruebas de significación y experimentos de exclusión de una fuente, que la literatura inglesa denomina \textit{Leave-One-Font-Out}, para cuantificar la capacidad de generalización a tipografías que el entrenamiento no cubre.
	
	\item Desarrollar la aplicación web para la digitalización, visualización en formato facsimilar y gestión de los documentos que procesa el sistema, e incluir la integración del modelo de reconocimiento y el diseño de roles para la revisión de transcripciones.
\end{enumerate}

La propuesta de solución se articula en cuatro componentes. El primero es un generador de datos de entrenamiento sintético que, a partir de obras literarias en español y de fuentes tipográficas de diversidad estilística, produce pares imagen--transcripción con una distribución de frecuencias de caracteres que corrige el desequilibrio natural del idioma. El segundo es un \textit{pipeline} de preprocesamiento que recibe documentos digitalizados y produce líneas de texto listas para el reconocimiento, con binarización adaptativa y normalización de tamaño que el sistema aplica en sucesión. El tercer componente, eje central del trabajo, es un modelo de reconocimiento basado en la arquitectura CRNN, que combina una red convolucional para la extracción de características con una red recurrente bidireccional para el modelado de secuencias, y que la función de pérdida CTC entrena; la inferencia decodifica mediante búsqueda en haz, que un modelo de lenguaje de cinco gramas en español guía durante la generación de la secuencia. El cuarto componente es la aplicación web, que integra el sistema y ofrece un visor de doble panel con el facsimilar original y la transcripción OCR, junto con funcionalidades de búsqueda, gestión y revisión de documentos con roles diferenciados.

La estructura de este documento es la siguiente. El capítulo~\ref{chap:estado-arte} introduce los fundamentos conceptuales del sistema: el flujo de procesamiento de un sistema OCR, las técnicas de preprocesamiento de imágenes de documentos, las propiedades estadísticas del lenguaje natural que justifican el diseño del conjunto de datos, y los fundamentos de las redes neuronales y el aprendizaje supervisado que el modelo de reconocimiento de caracteres requiere para su comprensión. El capítulo~\ref{chap:corpus} expone el análisis de los conjuntos de datos públicos disponibles, los fundamentos de la decisión de generar un \textit{corpus} propio, el diseño del sistema de generación y la validación estadística del \textit{corpus} que el procedimiento produce. El capítulo~\ref{chap:preprocesamiento} describe el \textit{pipeline} de preprocesamiento de documentos, que comprende la binarización adaptativa, la corrección de inclinación y la segmentación de líneas de texto. El capítulo~\ref{chap:modelo} presenta el diseño y entrenamiento del modelo de reconocimiento de secuencias, el protocolo de evaluación estadística que culmina con la medición sobre un conjunto de prueba independiente, y el mecanismo de decodificación por búsqueda en haz con modelo de lenguaje. El capítulo~\ref{chap:web} describe el desarrollo de la aplicación web, su arquitectura, las funcionalidades de digitalización y gestión documental, y el sistema de roles para la revisión de transcripciones. Por último, el documento recoge las conclusiones del trabajo, las recomendaciones para su extensión futura, las referencias bibliográficas y los anexos.

El desarrollo del proyecto contó con apoyo de herramientas de inteligencia artificial generativa, en sus versiones de acceso gratuito, para tareas mecánicas y repetitivas. Claude y Gemini intervinieron en la redacción del informe para la aplicación uniforme de las reglas de estilo a lo largo del documento, la detección de inconsistencias entre capítulos, y la corrección ortográfica y de puntuación. Perplexity y Gemini sirvieron como motor de búsqueda para localizar las referencias bibliográficas que el trabajo cita, verificar los datos de publicación y obtener explicaciones de la documentación técnica de bibliotecas externas. En el plano del código, las mismas herramientas asistieron en la generación de fragmentos repetitivos de la aplicación web —vistas, formularios y plantillas HTML que siguen los patrones estándar de Django—, en la adaptación de ejemplos de la documentación oficial a los tipos y a las estructuras de datos del proyecto, y en la escritura del código de acoplamiento entre el modelo de reconocimiento y el resto de la aplicación. Asimismo, contribuyeron a la confección de múltiples \textit{scripts} desechables para el análisis exploratorio del \textit{corpus} y la visualización gráfica de los resultados experimentales —histogramas de frecuencias de caracteres, distribuciones de longitudes de etiqueta, curvas de entrenamiento, comparativas entre configuraciones y paneles de diagnóstico del \textit{pipeline}—, a la generación de tablas \LaTeX{} a partir de los datos numéricos del proyecto, a la conversión entre formatos de serialización, y a la escritura de utilidades menores de inspección y validación del \textit{corpus} durante el desarrollo.
